\documentclass[10pt]{article}
\usepackage[margin=0.75in]{geometry}
\usepackage{amsmath, amssymb, amsthm}
\usepackage{microtype}
\usepackage{enumitem}
\usepackage{fancyhdr}
\usepackage{titlesec}
\usepackage{tcolorbox}
\usepackage{booktabs}

\linespread{1.08}
\setlength{\parskip}{0.4ex plus 0.2ex minus 0.1ex}

\titleformat{\section}{\large\bfseries}{\thesection.}{0.5em}{}
\titleformat{\subsection}{\normalsize\bfseries}{\thesubsection}{0.5em}{}
\titlespacing*{\section}{0pt}{1.5ex}{1ex}
\titlespacing*{\subsection}{0pt}{1ex}{0.5ex}

\setlist{itemsep=1pt, topsep=3pt, parsep=1pt, leftmargin=1.5em}

\pagestyle{fancy}
\fancyhf{}
\fancyhead[L]{\small EECS 127/227AT Discussion 1}
\fancyhead[R]{\small Spring 2026}
\fancyfoot[C]{\small\thepage}
\renewcommand{\headrulewidth}{0.4pt}

\newcommand{\RR}{\mathbb{R}}
\newcommand{\NN}{\mathcal{N}}
\newcommand{\Range}{\mathcal{R}}
\DeclareMathOperator{\rank}{rank}
\DeclareMathOperator{\trace}{tr}

\begin{document}

\begin{center}
{\Large\bfseries EECS 127/227AT \quad Optimization Models in Engineering \quad UC Berkeley \quad Spring 2026}\\[6pt]
{\LARGE\bfseries Discussion 1}
\end{center}

\section{Invertibility of $A^\top A$}

In this problem, we show that if the matrix $A \in \RR^{m \times n}$ has a full column rank, then the matrix $A^\top A$ is invertible.

\begin{enumerate}[label=(\alph*)]
\item Show that if a vector $\vec{x}$ is in the null space of $A$ then $\vec{x}$ is in the null space of $A^\top A$.

\textbf{Solution:}
\begin{align}
\vec{x} \in \NN(A) &\iff A\vec{x} = \vec{0} \\
&\implies A^\top A\vec{x} = \vec{0} \\
&\iff \vec{x} \in \NN\!\left(A^\top A\right)
\end{align}
Where line 2 follows by multiplying both sides of $A\vec{x} = 0$ by $A^\top$.

\item Conversely, show that if $\vec{x}$ is in the null space of $A^\top A$ then $\vec{x}$ is in the null space of $A$.

\textbf{Solution:}
\begin{align}
\vec{x} \in \NN\!\left(A^\top A\right) &\iff A^\top A\vec{x} = \vec{0} \\
&\implies \vec{x}^\top A^\top A\vec{x} = 0 \\
&\implies (A\vec{x})^\top A\vec{x} = 0 \\
&\implies \|A\vec{x}\|_2^2 = 0 \\
&\implies A\vec{x} = \vec{0} \\
&\implies \vec{x} \in \NN(A)
\end{align}
Where line 5 follows by multiplying both sides of $A^\top A\vec{x} = \vec{0}$ by $\vec{x}^\top$ and line 8 follows from the properties of norms.

\item Given that matrix $A$ has a full column rank, what can you say about its null space? What does this imply about the null space and invertibility of the matrix $A^\top A$?

\textbf{Solution:} $\NN(A) = \{\vec{0}\}$. From the previous parts, we have shown that $\NN(A) = \NN\!\left(A^\top A\right)$ then $\NN\!\left(A^\top A\right) = \{\vec{0}\}$ and thus $A^\top A$ is invertible.

\end{enumerate}

\newpage

\section{Eigenvalues}

Let $A \in \RR^{n \times n}$ have the eigendecomposition $P\Lambda P^{-1}$ where $\Lambda \in \RR^{n \times n}$ is a diagonal matrix with entries consisting of the eigenvalues $\lambda_1, \lambda_2, \ldots, \lambda_n$ and $P \in \RR^{n \times n}$ is an invertible matrix. Note that this is equivalent to stating that $A$ is diagonalizable via the transformation,
\begin{equation}
P^{-1}AP = \Lambda. \tag{10}
\end{equation}

\begin{enumerate}[label=(\alph*)]
\item Show that $A^m = P\Lambda^m P^{-1}$, for integer $m \geq 1$.

\textbf{Solution:}
\begin{align}
A^m &= (P\Lambda P^{-1})(P\Lambda P^{-1})\cdots(P\Lambda P^{-1}) \quad m \text{ times} \tag{11} \\
&= P\Lambda(P^{-1}P)\Lambda(P^{-1}P)\cdots\Lambda(P^{-1}P)\Lambda P^{-1} \tag{12} \\
&= P\Lambda^m P^{-1}. \tag{13}
\end{align}
The last equality follows from the repeated application of the identity $P^{-1}P = I$.

\item Show that determinant of $A$ is the product of its eigenvalues, i.e.
\begin{equation}
\det(A) = \prod_{i=1}^{n} \lambda_i. \tag{14}
\end{equation}

\textit{HINT: We have the identity} $\det(XY) = \det(X)\det(Y)$.

\textbf{Solution:} Write down eigendecomposition of $A$ and use properties of determinant given in the hint.
\begin{align}
\det(A) &= \det\!\left(P\Lambda P^{-1}\right) \tag{15} \\
&= \det(P)\det(\Lambda)\det\!\left(P^{-1}\right) \tag{16} \\
&= \det\!\left(PP^{-1}\right)\det(\Lambda) \tag{17} \\
&= \det(\Lambda) \tag{18} \\
&= \prod_{i=1}^{n} \lambda_i \tag{19}
\end{align}

\end{enumerate}

\vfill
\noindent\textcopyright{} UCB EECS 127/227AT, Spring 2026. All Rights Reserved. This may not be publicly shared without explicit permission.

\end{document}
